\section{Results}
\label{sec:results}

Paraphrasing had a minor impact on the accuracy of workflow generation and affected both
approaches to a similar degree. The MCP approach was consistently more accurate and reliable.
Workflow graphs that produced acceptable results but not identical to the correct graph are
labeled ``lenient'' in the results table (Table~\ref{tab:results}). MCP retrieval resulted in a
significantly larger context and token usage compared to the shortlist of 5~tools from RAG.
Quantitative analysis shows that all approved workflow graphs create acceptable results when run
with real geoprocesses.

% Results table of the extended abstract; numbers in results.csv.
% The original merged header cells ("Lenient accuracy" over three columns) are not allowed by
% the guidelines, so every column carries its full name.
\begin{table}[htbp]
  \caption{Results of MCP and RAG trial runs}
  \label{tab:results}
  \tablefont
  \begin{tabularx}{\textwidth}{|l|*{8}{>{\centering\arraybackslash}X|}}
    \hline
    \rowcolor{tablehead}
    \textbf{Variant}
      & \textbf{Lenient seed} & \textbf{Lenient para\-phrases} & \textbf{Lenient all}
      & \textbf{Strict seed} & \textbf{Strict para\-phrases} & \textbf{Strict all}
      & \textbf{Input tokens/ req} & \textbf{Output tokens/ req} \\
    \hline
    \csvreader[late after line=\\\hline]{tables/results.csv}{
      variant=\variant,
      lenient_seed=\lenientseed, lenient_paraphrases=\lenientparaphrases, lenient_all=\lenientall,
      strict_seed=\strictseed, strict_paraphrases=\strictparaphrases, strict_all=\strictall,
      tokens_input=\tokensinput, tokens_output=\tokensoutput
    }{%
      \variant
      & \lenientseed\% & \lenientparaphrases\% & \lenientall\%
      & \strictseed\% & \strictparaphrases\% & \strictall\%
      & \tokensinput & \tokensoutput}
  \end{tabularx}
  \tabsource{\ownwork}
\end{table}


MCP provided more accurate and reliable performance with lenient and strict scores of 95.3~\%
and 92.3~\% respectively. Unlike the RAG results, the MCP produced few graphs that were only
leniently correct. Instead, the vast majority of them only made the correct graph workflow.

On the other hand, RAG provided graphs that were leniently accurate 85.9~\% of the time but
strictly accurate only 27.1~\% of the time. This was due to the frequent inclusion of an extra
buffer node at the end of the workflow, in addition to the correct interval buffer node.

When considering only the lenient results, there was not that much of a difference between the
two approaches. The difference was roughly 10~\% poorer for the RAG retrieval method. Paraphrasing
had a minor impact on the performance of both RAG and MCP retrieval. For MCP, the impact was only
observed when scoring against the strict criteria. Meanwhile, RAG showed negative effects only
for the lenient criteria. In both cases, the observed effect was small, at less than
5~percentage points.

Where MCP clearly had a drawback was in token usage. Each time the retrieval node called the list
of tools, this was added as context to the LLM request. This resulted in an average of 191K
tokens of input per request for MCP compared to an average of 88K for RAG. Output tokens remained
roughly the same, with MCP having a slightly lower token cost.

\begin{figure}[htbp]
  \caption{An example of correct output visualized on a map, dark colored regions indicate
    ``cool spots''}
  \label{fig:map}
  \includegraphics[width=\textwidth, alt={Map of the south-western part of Münster. Dark
    shading covers the areas within the walking-distance buffers around large parks and forests;
    the remaining unshaded patches, mostly between the districts, lie outside the 300 m
    limit.}]{04_result-map}
  \figsource{\ownwork; map data \textcopyright{} OpenStreetMap contributors,
    \url{https://www.openstreetmap.org/copyright}}
\end{figure}

\begin{figure}[htbp]
  \caption{Workflow output visualized in the user interface}
  \label{fig:ui}
  \includegraphics[width=\textwidth, alt={Screenshot of the execution graph view in the web
    interface: five nodes (city-boundary, osm-green-areas, min-area-filter, buffer and
    multi-ring-buffer) connected by four edges and marked as valid, with buttons to delete the
    graph, show results and run it.}]{05_result-ui}
  \figsource{\ownwork}
\end{figure}

The resulting graphs were validated by executing the workflows with the actual services of the
geoserver (Figures~\ref{fig:map} and~\ref{fig:ui}). The exact quantities that resulted depended
on arbitrary decisions made by the agent, for example the number of parks within the city
boundary that meet the criteria or the minimum size of the park. For this reason, the results
were validated by checking that they did not vary significantly from a known correct solution.
By these metrics, all graphs deemed correct, both lenient and strict, did not differ to such a
degree that they would be considered incorrect.
